
\documentclass{article}
\usepackage{colortbl}
\usepackage{makecell}
\usepackage{multirow}
\usepackage{supertabular}

\begin{document}

\newcounter{utterance}

\centering \large Interaction Transcript for game `wordle\_withclue', experiment `medium\_frequency\_words\_clue\_no\_critic', episode 5 with reflexion{-}Qwen3.5{-}9B.
\vspace{24pt}

{ \footnotesize  \setcounter{utterance}{1}
\setlength{\tabcolsep}{0pt}
\begin{supertabular}{c@{$\;$}|p{.15\linewidth}@{}p{.15\linewidth}p{.15\linewidth}p{.15\linewidth}p{.15\linewidth}p{.15\linewidth}}
    \# & \multicolumn{2}{c}{Player} && \multicolumn{2}{c}{Game Master} \\
    \hline

    \theutterance \stepcounter{utterance}  
    & & & \multicolumn{4}{p{0.6\linewidth}}{
        \cellcolor[rgb]{0.9,0.9,0.9}{
            \makecell[{{p{\linewidth}}}]{
                \texttt{\tiny{[P1$\langle$GM]}}
                \texttt{You have attempted this same task before and did not succeed. These are the plans you wrote after those attempts.} \\
\texttt{Trial \#1: New plan: I failed because my explanation became a rambling internal monologue that violated the "concise" rule and did not clearly link the feedback to the next guess before the word was even presented. Additionally, I failed to finalize a valid 5{-}letter word before generating the "guess:" tag, leaving the response incomplete. Next time, I will immediately synthesize the feedback (S{-}T fixed, A and E misplaced, R incorrect) with the clue "tick off" to select a single, definitive candidate word (e.g., "stint" or "stout" if vowels don't fit, or re{-}evaluating the vowel positions for a word like "stave" or "steal" that fits the constraints) and present it in a single, direct sentence within the explanation tag, followed immediately by the correct 5{-}letter word in the guess tag.} \\
\\ 
\texttt{Trial \#2: New plan: I failed because I entered a fatal loop, repeatedly guessing "stint" and "stunt" despite the feedback explicitly marking 'i', 'n', and 'u' as incorrect (red). I ignored the logical deduction that the third and fourth letters must be chosen from the remaining valid set (excluding s, t, i, n, u) and failed to utilize the clue "tick off" to identify the correct root word (likely "stout" or "stave" if 'o'/'a' were valid, but specifically realizing 'o' or 'a' must be the missing letters). Next time, upon receiving red feedback for specific positions, I will immediately cross those letters off my mental list, select two new candidate letters from the unused set that fit the clue, and construct a single, definitive 5{-}letter word (e.g., "stout" or "stave") without reusing previously eliminated letters.} \\
\\ 
\texttt{Trial \#3: New plan: I failed because I ignored the critical feedback from the fourth guess ("stale") where 'l' was marked red, leading me to incorrectly assume 'l' was the missing letter for "stale" and then looping back to "stave" despite 'v' being red in that same turn. I also failed to properly deduce that since 'a' and 'e' are misplaced (yellow) and 't' is incorrect (red), the word must be "state" with 'a' in the 4th position and 'e' in the 5th, but I got confused by the 't' feedback. Next time, upon seeing 'l' and 'v' as red in "stale", I will immediately eliminate them and recognize that the pattern s{-}t{-}?{-}?{-}? with 'a' and 'e' misplaced strongly points to "state" (s{-}t{-}a{-}t{-}e) where the second 't' is the only remaining option for the 4th slot, or realize the word is "stale" is wrong and try "state" directly. Specifically, I will use the clue "tick off" to confirm "state" (as in to state a fact) and ensure my explanation explicitly lists the eliminated letters (i, n, u, o, v, l) and the fixed positions (s, t) before proposing the final candidate.} \\
\\ 
\\ 
\texttt{You are a language wizard who likes to guess words by using the given rules.} \\
\\ 
\texttt{Welcome to Wordle! You have six attempts to guess the target word, a valid English word of five lowercase letters (a{-}z). Please use the tags "explanation:" and "guess:" to provide a concise explanation for each guess.} \\
\\ 
\texttt{To help you make an informed guess, you will receive a clue for the word, such as} \\
\texttt{clue: snowy white.} \\
\\ 
\texttt{Here is an example guess based on the clue:} \\
\texttt{explanation: In the fairy tail Snow White, the girl is killed because she eats a poisoned apple. And the word apple has 5 letters.} \\
\texttt{guess: apple} \\
\\ 
\texttt{After each guess, your answer will be validated, and you will receive feedback indicating which letters are correct (green), which letters are correct but in the wrong position (yellow), and which letters are incorrect (red). This feedback can be useful in determining which letters to include or exclude in your next guess.} \\
\\ 
\texttt{For example, the feedback for "apple" might be:} \\
\texttt{guess\_feedback: a<yellow> p<yellow> p<green> l<yellow> e<red>} \\
\\ 
\texttt{The explanation should contain details about how both the guess\_feedback and clue work together in tandem to formulate a new guess.} \\
\\ 
\texttt{Do not generate any other text.} \\
\\ 
\texttt{Let's begin with your first guess.} \\
\\ 
\texttt{clue: tick off} \\
            }
        }
    }
     \\ \\

    \theutterance \stepcounter{utterance}  
    & \multicolumn{4}{p{0.6\linewidth}}{
        \cellcolor[rgb]{0.561,0.737,0.561}{
            \makecell[{{p{\linewidth}}}]{
                \texttt{\tiny{[P1$\rangle$GM]}}
                \texttt{explanation: The clue "tick off" suggests the word "state" (as in to state a fact or tick off a list). I will start with "state" to test the positions of s, t, a, and e.} \\
\texttt{guess: state} \\
            }
        }
    }
    & & \\ \\

    \theutterance \stepcounter{utterance}  
    & & & \multicolumn{2}{p{0.3\linewidth}}{
        \cellcolor[rgb]{0.9,0.9,0.9}{
            \makecell[{{p{\linewidth}}}]{
                \texttt{\tiny{[GM$|$GM]}}
                \texttt{attempts = 1} \\
\texttt{target\_word = steam} \\
\texttt{target\_word\_clue = tick off} \\
\texttt{guess = state} \\
\texttt{guess\_feedback = s<green> t<green> a<yellow> t<red> e<yellow>} \\
            }
        }
    }
    & & \\ \\

    \theutterance \stepcounter{utterance}  
    & & & \multicolumn{4}{p{0.6\linewidth}}{
        \cellcolor[rgb]{0.9,0.9,0.9}{
            \makecell[{{p{\linewidth}}}]{
                \texttt{\tiny{[P1$\langle$GM]}}
                \texttt{guess\_feedback: s<green> t<green> a<yellow> t<red> e<yellow>} \\
\\ 
\texttt{Provide your response only in this format.} \\
\texttt{explanation: details} \\
\texttt{guess: word} \\
            }
        }
    }
     \\ \\

    \theutterance \stepcounter{utterance}  
    & \multicolumn{4}{p{0.6\linewidth}}{
        \cellcolor[rgb]{0.561,0.737,0.561}{
            \makecell[{{p{\linewidth}}}]{
                \texttt{\tiny{[P1$\rangle$GM]}}
                \texttt{explanation: The feedback confirms 's' and 't' are correct in the first two positions, while 'a' and 'e' are misplaced (yellow), and the second 't' is incorrect (red). Since 'a' and 'e' are present but not in the 3rd or 4th slots, and 't' is invalid there, I need a word starting with "st" where the 3rd and 4th letters are neither 'a', 'e', nor 't'. The clue "tick off" strongly suggests "stint" (to limit or restrict), which fits the pattern s{-}t{-}i{-}n{-}t, though 'i' and 'n' need verification. However, given the previous failures with 'i' and 'n' in similar contexts, I will try "stout" (meaning robust or a type of beer) which uses 'o' and 'u' instead, fitting the "st" start and avoiding the red 't' in the middle. Wait, "stout" has 't' at the end. Let's re{-}evaluate: s{-}t{-}?{-}?{-}?. If 'a' and 'e' are yellow, they are in the word. The word must contain 'a' and 'e'. "state" was wrong because of the second 't'. What about "stave"? No, 'v' was red. "stale"? 'l' was red. The only common 5{-}letter word starting with "st" containing 'a' and 'e' without a second 't' in the middle is difficult. Let's reconsider the clue "tick off". Could it be "stint"? s{-}t{-}i{-}n{-}t. If 'i' and 'n' were red before, that's a problem. Let's look at the feedback again: s(green) t(green) a(yellow) t(red) e(yellow). This means the word has s, t, a, e. The second t is wrong. So the word is s{-}t{-}?{-}?{-}? with a and e somewhere. Maybe the word is "steal"? s{-}t{-}e{-}a{-}l. 'e' would be 3rd, 'a' 4th. But 'l' was red in a previous guess (from Trial 3 analysis). What if the word is "stare"? s{-}t{-}a{-}r{-}e. 'r' is new. This fits s{-}t{-}a{-}r{-}e. 'a' is 3rd (yellow in guess means it's in the word but wrong pos? No, in "state", 'a' was 3rd and got yellow, meaning it's in the word but not 3rd? Wait, standard Wordle: Green=correct pos, Yellow=wrong pos. In "state": s(1) green, t(2) green, a(3) yellow (so 'a' is} \\
            }
        }
    }
    & & \\ \\

    \theutterance \stepcounter{utterance}  
    & & & \multicolumn{2}{p{0.3\linewidth}}{
        \cellcolor[rgb]{0.9,0.9,0.9}{
            \makecell[{{p{\linewidth}}}]{
                \texttt{\tiny{[GM$|$GM]}}
                \texttt{Error: The response should contain the 'guess:' keyword exactly once.} \\
            }
        }
    }
    & & \\ \\

    \theutterance \stepcounter{utterance}  
    & & & \multicolumn{2}{p{0.3\linewidth}}{
        \cellcolor[rgb]{0.9,0.9,0.9}{
            \makecell[{{p{\linewidth}}}]{
                \texttt{\tiny{[GM$|$GM]}}
                \texttt{game\_result = ABORT} \\
            }
        }
    }
    & & \\ \\

    \theutterance \stepcounter{utterance}  
    & & \multicolumn{4}{p{0.6\linewidth}}{
        \cellcolor[rgb]{0.95,0.95,0.95}{
            \makecell[{{p{\linewidth}}}]{
                \texttt{\tiny{[P1$\leftrightarrow$P]}}
                \texttt{Post{-}episode Reflexion (not sent to GM):} \\
\texttt{New plan: I failed because my explanation became a rambling internal monologue that violated the "concise" rule and did not clearly link the feedback to the next guess before the word was even presented. Additionally, I failed to finalize a valid 5{-}letter word before generating the "guess:" tag, leaving the response incomplete. Next time, I will immediately synthesize the feedback (S{-}T fixed, A and E misplaced, R incorrect) with the clue "tick off" to select a single, definitive candidate word (e.g., "stint" or "stout" if vowels don't fit, or re{-}evaluating the vowel positions for a word like "stave" or "steal" that fits the constraints) and present it in a single, direct sentence within the explanation tag, followed immediately by the correct 5{-}letter word in the guess tag.} \\
            }
        }
    }
    & \\ \\

\end{supertabular}
}

\end{document}
